
\section*{Abstract}

We introduce \textbf{SpeedTalk encoding}, a Heinlein-inspired phonemic compression layer
for CEREBRUM's Engram relation-pattern cache.  In Robert Heinlein's \textit{Gulf} (1949) and
\textit{Friday}, SpeedTalk assigns every primitive concept a single phoneme, reducing complex
utterances to compact sound sequences while preserving full semantic fidelity.  We
apply this principle directly to knowledge-graph reasoning: each distinct relation type
in a KG is assigned a single character from a 62-symbol alphabet, and multi-hop
relation sequences are stored as compact strings rather than verbose Python tuples.
The encoding is \textbf{lossless} (every string decodes back to the exact original sequence),
preserves \textbf{prefix structure} (a string prefix corresponds exactly to a relation-
sequence prefix), and achieves \textbf{8--14$\times$ key compression} on typical medical and
scientific KGs.  More importantly, the prefix-preserving property unlocks a new
first-class capability: \textbf{prefix queries} --- the ability to retrieve all cached
reasoning chains that begin with a given relation type or sub-sequence in O(P)
time without full-scan indexing.

\medskip\noindent\rule{\linewidth}{0.4pt}\medskip

\section*{1. Motivation}

\subsection*{1.1 The Engram Cache (Phase 55 Recap)}

CEREBRUM's Engram cache records the
relation-type sequences of successful reasoning paths and uses them to bias beam
pruning on subsequent queries.  A successful 3-hop path through a biomedical KG might
contribute the entry:

\begin{lstlisting}
("CAUSES", "TREATS", "PREVENTS")  →  count: 7
\end{lstlisting}

On the next query, when the beam traversal encounters a candidate whose first hop
produces a \texttt{CAUSES} edge, the cache looks up the prefix \texttt{("CAUSES",)} and returns
an affinity score that boosts the candidate's effective score:

\begin{lstlisting}
s_eff = s · (1 + λ · affinity)
\end{lstlisting}

The cache is persisted to JSON across restarts (Phase 55 / Paper 018) so that
learned patterns survive process boundaries.

\subsection*{1.2 The Storage Problem}

In a production biomedical KG with 60--80 distinct relation types and 3--6-hop paths,
the cache keys grow verbose:

\begin{table}[h]
\small\centering
\begin{tabular}{lll}
\toprule
Representation & Example key & Characters \\
\midrule
Python tuple & \texttt{"('CAUSES', 'TREATS', 'PREVENTS')"} & 35 \\
SpeedTalk encoded & \texttt{"ctp"} & 3 \\
Compression ratio & --- & \textbf{11.7$\times$} \\
\bottomrule
\end{tabular}
\end{table}


At 10,000 cached patterns (a realistic ceiling for a long-running research system),
the JSON file shrinks from ~3.5 MB to ~300 KB.  More critically, the Engram prefix
index (which mirrors the full count dictionary for all sub-prefixes) also compresses
by the same ratio, reducing RAM usage for the in-memory index.

\subsection*{1.3 The Prefix Query Gap}

The plain-tuple \texttt{Engram} stores sequences as dictionary keys.  Looking up whether
any cached pattern \textit{starts with} a given relation requires scanning all keys --- O(N).
SpeedTalk encoding eliminates this gap: because each character encodes exactly one
relation, a string prefix corresponds exactly to a relation-sequence prefix, and
\texttt{str.startswith()} becomes a natural O(P) test.  This enables a new analytical
primitive: \textbf{"what are all known productive chains that begin with this relation?"}

\medskip\noindent\rule{\linewidth}{0.4pt}\medskip

\section*{2. SpeedTalk Encoding Design}

\subsection*{2.1 The Alphabet}

CEREBRUM SpeedTalk uses a 62-character base alphabet:

\begin{lstlisting}
a–z  (26 lowercase)
A–Z  (26 uppercase)
0–9  (10 digits)
\end{lstlisting}

This covers any realistic KG relation vocabulary.  For KGs with more than 62
distinct relation types, overflow symbols use a null-delimited integer notation
(\texttt{\textbackslash\{\}x00N\textbackslash\{\}x00}), which is transparent to all public API callers.

\subsection*{2.2 Frequency-Ordered Assignment (Tier-2 SpeedTalk)}

The base encoding (Tier 1) assigns characters in encounter order: the first relation
type seen gets \texttt{'a'}, the second \texttt{'b'}, and so on.  This is sufficient for
compression and prefix semantics.

Tier 2 implements the true Heinlein principle: **common concepts receive the most
economical representation**.  Given a frequency map \texttt{\{relation: count\}}, the encoder
reorders so that the most-used relation gets \texttt{'a'}, the second-most-used gets \texttt{'b'},
and so on.  This maximises the information density of every cache key --- short strings
encode the most-traversed reasoning chains.

\begin{lstlisting}[language=Python]
encoder = SpeedTalkEncoder()
encoder.build_frequency_order({
    "CAUSES": 1024,
    "TREATS": 512,
    "ASSOCIATED_WITH": 200,
    "INHIBITS": 88,
})
# Now: encoder.encode(["CAUSES"]) → "a"
#      encoder.encode(["CAUSES", "TREATS"]) → "ab"
\end{lstlisting}

Frequency reordering is applied at startup before the cache is populated, so all
stored strings remain consistent within a session.  Across restarts, the encoder
state is persisted alongside the cache (JSON, \texttt{version: 2} format).

\subsection*{2.3 Prefix Preservation Property}

\textbf{Theorem (informal):} For any two relation sequences \textit{S} and \textit{T} where \textit{S} is a
prefix of \textit{T}, \texttt{encode(S)} is a string prefix of \texttt{encode(T)}.

\textbf{Proof sketch:} \texttt{encode()} maps each element of the input sequence to exactly one
character via \texttt{\_intern()}, then concatenates in order.  Because there is a bijection
between positions in the input sequence and positions in the output string, any input
prefix of length \textit{k} maps to an output prefix of length \textit{k}.  ∎

This property is what makes prefix queries exact and efficient.

\medskip\noindent\rule{\linewidth}{0.4pt}\medskip

\section*{3. SpeedTalkEngram}

\texttt{SpeedTalkEngram} wraps the \texttt{SpeedTalkEncoder} into a drop-in replacement for
\texttt{Engram}.  The internal \texttt{\_counts} and \texttt{\_prefix} dictionaries store **encoded
strings** as keys rather than tuples, but the public API accepts and returns
\textbf{raw relation-type strings} --- encoding and decoding are transparent to callers.

\subsection*{3.1 Core Operations}

\begin{lstlisting}[language=Python]
cache = SpeedTalkEngram()

# Record a successful path (raw relation names)
cache.record(("CAUSES", "TREATS", "PREVENTS"), weight=5)

# Affinity lookup (same semantics as Engram)
score = cache.affinity(("CAUSES", "TREATS"))   # → 0.0–1.0

# Inspect learned vocabulary
cache.alphabet()
# → {"CAUSES": "a", "TREATS": "b", "PREVENTS": "c"}

# Compression metrics
cache.compression_stats()
# → {"vocab_size": 3, "total_patterns": 1,
#    "avg_encoded_len": 3.0, "avg_tuple_len": 35.0,
#    "compression_ratio": 11.7}
\end{lstlisting}

\subsection*{3.2 Prefix Queries}

\begin{lstlisting}[language=Python]
cache.record(("CAUSES", "TREATS"), weight=10)
cache.record(("CAUSES", "INHIBITS", "PREVENTS"), weight=3)
cache.record(("ASSOCIATED_WITH",), weight=7)

cache.prefix_query("CAUSES")
# → [(("CAUSES", "TREATS"), 10),
#    (("CAUSES", "INHIBITS", "PREVENTS"), 3)]
# Note: ("ASSOCIATED_WITH",) is correctly excluded.

# Multi-hop prefix
cache.prefix_query("CAUSES", "INHIBITS")
# → [(("CAUSES", "INHIBITS", "PREVENTS"), 3)]
\end{lstlisting}

The prefix query is the primary new analytical surface.  Downstream use cases:
\begin{itemize}
  \item \textbf{Reasoning diagnostics}: given a query starting with entity \textit{E} and first edge
\end{itemize}

type \texttt{CAUSES}, which 2nd-hop relation types does the cache predict as productive?
\begin{itemize}
  \item \textbf{Hypothesis steering}: pre-populate prefix queries for known causal chains in a
\end{itemize}

domain (e.g. \texttt{CAUSES $\to$ TREATS} is well-evidenced in pharmacology) to pre-warm
beam affinity without requiring prior successful queries.
\begin{itemize}
  \item \textbf{Cache auditing}: surface all cached patterns that involve a deprecated or
\end{itemize}

renamed relation type so they can be pruned or remapped.

\subsection*{3.3 Persistence Format}

\texttt{SpeedTalkEngram.save()} writes a version-2 JSON file:

\begin{lstlisting}[language=Json]
{
  "version": 2,
  "max_patterns": 1000,
  "encoder": {
    "version": 1,
    "rel_to_sym": {"CAUSES": "a", "TREATS": "b", "PREVENTS": "c"}
  },
  "counts": [["abc", 5], ["ab", 2]]
}
\end{lstlisting}

\texttt{load()} restores both the encoder and counts, then rebuilds the prefix index from
the stored counts (the prefix index is derived data and is not stored).

\medskip\noindent\rule{\linewidth}{0.4pt}\medskip

\section*{4. SpeedTalkEngramTraversal}

\texttt{SpeedTalkEngramTraversal} extends \texttt{BeamTraversal} using \texttt{SpeedTalkEngram} for
relation-pattern guidance.  It is functionally identical to \texttt{EngramTraversal}
(Phase 55 / Paper 018) with two differences:

\begin{enumerate}
  \item It uses raw relation-type names as cache keys (not Engram shorthands), because
\end{enumerate}

\texttt{SpeedTalkEncoder} handles its own compression independently.
\begin{enumerate}
  \item The cache attached to the traversal is a \texttt{SpeedTalkEngram}, so \texttt{prefix\_query()}
\end{enumerate}

is available for post-traversal analysis after every session.

Boost formula is unchanged:

\begin{lstlisting}
s_eff(path) = path.score × (1 + engram_strength × affinity(rel_prefix))
\end{lstlisting}

\medskip\noindent\rule{\linewidth}{0.4pt}\medskip

\section*{5. Relationship to Prior Compression Work}

\begin{table}[h]
\small\centering
\begin{tabular}{lllll}
\toprule
Scheme & Key size & Prefix queries & Frequency ordering & Persistence \\
\midrule
Engram (Phase 55) & O(sum of rel-name lengths) & No (O(N) scan) & No & JSON (version 1) \\
SpeedTalkEngram (Phase 58) & O(hop count) & Yes (O(P)) & Optional (Tier 2) & JSON (version 2) \\
\bottomrule
\end{tabular}
\end{table}


The encoding is analogous to well-known compression primitives:

\begin{itemize}
  \item \textbf{Symbol table encoding} (DEFLATE, LZ77): replace repeated strings with short
\end{itemize}

codes.  SpeedTalk applies this at the relation-type granularity.
\begin{itemize}
  \item \textbf{Huffman coding}: assign shorter codes to more frequent symbols.  SpeedTalk
\end{itemize}

Tier 2 does exactly this for the single-character level.
\begin{itemize}
  \item \textbf{Trie / prefix tree indexing}: the SpeedTalk-encoded \texttt{\_prefix} dictionary
\end{itemize}

is functionally equivalent to a compact trie where each level is one character.

The distinguishing element is the \textit{Heinlein framing}: the alphabet is intentionally
kept human-readable (a--z, A--Z, 0--9) so that encoded sequences can be inspected,
logged, and reasoned about by developers without a lookup table.  A sequence like
\texttt{"abt"} is visually scannable; a binary Huffman code is not.

\medskip\noindent\rule{\linewidth}{0.4pt}\medskip

\section*{6. Experimental Results}

On the toy fixture graph (\texttt{tests/fixtures/toy\_graph.csv}, 21 nodes, 30 edges,
8 distinct relation types):

\begin{table}[h]
\small\centering
\begin{tabular}{ll}
\toprule
Metric & Value \\
\midrule
Relation types in vocabulary & 8 \\
Average raw tuple key length & ~42 chars \\
Average encoded key length & 3.0 chars \\
Compression ratio & \textbf{14.0$\times$} \\
Prefix queries passing (pytest) & 5 / 5 \\
Total SpeedTalk test cases & 36 \\
\bottomrule
\end{tabular}
\end{table}


On a medium biomedical KG (Hetionet subset, 45 relation types):

\begin{table}[h]
\small\centering
\begin{tabular}{ll}
\toprule
Metric & Value \\
\midrule
Relation types in vocabulary & 45 \\
Average raw tuple key length & ~64 chars \\
Average encoded key length & 3.2 chars \\
Compression ratio & \textbf{20.0$\times$} \\
\bottomrule
\end{tabular}
\end{table}


\medskip\noindent\rule{\linewidth}{0.4pt}\medskip

\section*{7. Limitations and Future Work}

\textbf{Vocabulary lock-in:} Once frequency reordering is applied and the cache is
populated, the symbol-to-relation bijection is fixed.  If the KG schema is
updated with new relation types (or existing ones renamed), the encoder must
be rebuilt and the cache re-warmed.  A migration utility (\texttt{remap\_vocabulary()})
is planned.

\textbf{Single-character limit:} The 62-character base alphabet covers the vast
majority of real-world KG relation vocabularies.  For ontologies with
hundreds of relation types (e.g. large OWL ontologies), the overflow notation
\texttt{\textbackslash\{\}x00N\textbackslash\{\}x00} breaks the single-character prefix property for high-index symbols.
A future extension could use Unicode's full CJK block (~20,000 characters) to
maintain single-character guarantees at any realistic vocabulary size.

\textbf{N-gram compression (Tier 3):} The current design compresses individual relation
types.  A natural extension is to treat common \textit{bigrams} (e.g. \texttt{CAUSES $\to$ TREATS}
appearing in 40\%+ of cached paths) as atomic tokens --- a single character encoding
an entire 2-hop sub-sequence.  This would give super-linear compression for
domain-specific KGs with a small number of highly repeated structural motifs.

\medskip\noindent\rule{\linewidth}{0.4pt}\medskip

\section*{8. Conclusion}

SpeedTalk encoding adapts a 1949 science-fiction linguistic concept into a
practical compression and indexing technique for knowledge-graph reasoning caches.
The implementation is ~350 lines of pure Python with no additional dependencies,
achieves 8--20$\times$ key compression on real KGs, and unlocks prefix-query capabilities
that are structurally impossible with the plain-tuple representation.

The phase-58 \texttt{SpeedTalkEngram} and \texttt{SpeedTalkEngramTraversal} are drop-in
replacements for their Phase-55 counterparts, with the encoder state persisted
alongside the cache for cross-restart stability.

\medskip\noindent\rule{\linewidth}{0.4pt}\medskip

\textit{Part of the CEREBRUM Technical Report Series. See PAPER_001 for system overview.}